\documentclass{shnu-exam}

\examtitle{概率论与数理统计}
\examtype{A卷}
\semester{2025~2026 学年~~第~1~学期}
\examdate{2026年1月13日}
\examduration{90}

\begin{document}
\maketitle

\section{一、选择题（每小题2分，共16分）}

\question 设随机变量$X$和$Y$相互独立，且$X\sim N(1,4)$，$Y\sim N(2,9)$，则$Z=X-Y$服从的分布为
\choice{$N(-1,13)$}
\choice{$N(3,13)$}
\choice{$N(-1,5)$}
\choice{$N(3,5)$}

\question 下列关于大数定律的表述，正确的是
\choice{切比雪夫大数定律要求随机变量序列相互独立且同分布}
\choice{辛钦大数定律可用于判断样本均值的分布类型}
\choice{伯努利大数定律验证了"频率的稳定性"，是大数定律的重要应用}
\choice{大数定律表明，样本量越大，样本均值与总体均值的偏差一定越小}

\question 在假设检验中，若显著性水平$\alpha=0.05$，下列说法正确的是
\choice{拒绝原假设$H_0$的概率为0.05}
\choice{接受原假设$H_0$的概率为0.05}
\choice{原假设$H_0$成立时，拒绝$H_0$的概率为0.05}
\choice{备择假设$H_1$成立时，接受$H_0$的概率为0.05}

\section{二、填空题（每空3分，共24分）}

\question 设事件$A,B$满足$P(A)=0.4$，$P(B)=0.5$，$P(A|B)=0.6$，则$P(A\cup B)=$\fillin.

\question 设随机变量$X$服从区间$[0,2]$上的均匀分布，则$P(0.5<X<1.5)=$\fillin.

\question 设随机变量$X\sim B(10,0.2)$（二项分布），则$E(X)=$\fillin.

\section{三、计算题（本题10分）}

\problem[lines=15]{某实验室有三台检测设备A、B、C。在所有检测任务中，有50\%由设备A完成，30\%由设备B完成，剩下20\%由设备C完成。已知设备A的误判率为1\%，设备B的误判率为2\%，设备C的误判率为4\%。现从所有已经完成的检测任务中随机抽取一次，发现该次检测结果是误判。请问这次检测是由设备B完成的可能性有多大？}

\section{四、计算题（本题16分）}

\problem[space=8cm]{设二维随机变量$(X,Y)$的联合概率密度函数为
$f(x,y)=\begin{cases}4xy,&0<x<1,x^2<y<x\\0,&\text{其他}\end{cases}$
(1) 求$X$和$Y$的边缘密度函数；(2) $X$和$Y$是否独立？}

\end{document}
